\section{Introduction}
A fixed per-image channel-symbol budget limits the source description available for reconstruction. Feasibility depends on image-dependent source-code length and in-frame protection. A generative prior completes untransmitted information, constrained by the delivered description.

Digital coding and learned joint source--channel coding (JSCC) offer different solutions~\cite{deepjscc,swinjscc}; fitting progressive descriptions to channel resources is a classical problem~\cite{chande}. Ada-TokenCom combines prefix entropy coding, tail generation, and prefix/MCS selection under an average resource constraint~\cite{adatokencom}. TONIC addresses task-oriented token transmission~\cite{tonic}, and ARPC uses next-scale progressive coding~\cite{arpc}. We adopt visual autoregressive modeling (VAR), which reported state-of-the-art results in its original class-conditional ImageNet generation benchmarks~\cite{var}. Its complete coarse maps describe the full latent grid after resizing; later maps add residual refinements. This organization motivates transmitting a limited coarse-to-fine description for source-constrained reconstruction.

Whole-scale boundaries need not align with frame capacity: a complete scale adds many tokens, and its terminated code length varies across images. Within-scale endpoints use capacity left after the last feasible whole scale. Both ends construct matching coding probabilities from complete coarse maps. The received boundary tokens must then enter residual accumulation and the conditions for finer-scale prediction.

NeST-Com combines within-scale endpoints, independently terminated arithmetic streams, actual-length fitting, and prefix-constrained reconstruction. It uses VAR's known-token overwrite~\cite{var} to merge received boundary positions with predictions and update both residual latent and finer-scale context. Shared-model comparisons show 0.32--0.70~dB PSNR gains over whole-scale transmission at $N=1024$ and 0.39--0.79~dB at $N=512$. Controlled endpoint and receiver interventions examine the design, while six-system comparisons and separately calibrated reduced-budget results characterize complete-system quality tradeoffs.
